% GENERATED FILE. Edit canonical lessons, then run npm run generate:books.
% canonical-source-hash: fnv1a64:6d4d896aaf4a52a3
% canonical-lessons: FA-C502-gahvare, FA-C502-lalayi, FA-C502-pushak, FA-C502-kaleske, FA-C502-shir-dadan

\chapter{Babies}
\label{ch:fa-babies}
\hlchaptermodality{502}

\begin{chapteropening}
\textbf{By the end of this chapter:} \emph{I can talk about looking after a baby.}

Complete the last lesson of chapter 502: I can talk about looking after a baby.
\end{chapteropening}

\section[\texorpdfstring{\fa{گهواره} (gahvâre)}{گهواره (gahvâre)}]{\fa{گهواره} (gahvâre) — a cradle}

\label{lesson:FA-C502-gahvare}



\begin{quote}
Before the new one: say the Persian for twins, then the Persian for a child.
\end{quote}

\subsection*{You'll want to know: \fa{گهواره}}
\textbf{\fa{گهواره}} (\emph{gahvâre}) — \textquotedblleft{}a cradle\textquotedblright{}.

\begin{quote}
\fa{مادر} \fa{بچه} \fa{را} \fa{در} \fa{گهواره} \fa{گذاشت}.
\end{quote}

\begin{quote}
\emph{mâdar bachche râ dar gahvâre gozâsht.}
\end{quote}

\begin{quote}
The mother put the baby in the cradle.
\end{quote}

Plural \emph{gahvârehâ}.

\begin{cousinweb}
\emph{gahvâre} is inherited: EtymDB traces it to a Middle Persian word meaning 'cot, cradle' (EtymDB).
\end{cousinweb}

\begin{grammarlens}[title={What you've built}]
One more word for telling, explaining and describing at B1.
\end{grammarlens}

\subsection*{Your turn}
\begin{itemize}
\raggedright
  \item \emph{Say it:} \emph{gahvâre}
  \item \emph{Say it:} \emph{gahvâre}, once more
  \item \emph{Say it:} a sentence of your own with \emph{gahvâre}
  \item \emph{Recall:} say the Persian for twins, then the Persian for a child, then say \emph{gahvâre} again
\end{itemize}

\begin{tcolorbox}[breakable,colback=teal!4,colframe=teal!35!black,title={Before you move on}]
What does \fa{گهواره} mean? (\textquotedblleft{}A cradle\textquotedblright{}.) Say it once more.
\end{tcolorbox}
\section[\texorpdfstring{\fa{لالایی} (lâlâyi)}{لالایی (lâlâyi)}]{\fa{لالایی} (lâlâyi) — a lullaby}

\label{lesson:FA-C502-lalayi}



\begin{quote}
Before the new one: say the Persian for a child, then the Persian for a cradle.
\end{quote}

\subsection*{You'll want to know: \fa{لالایی}}
\textbf{\fa{لالایی}} (\emph{lâlâyi}) — \textquotedblleft{}a lullaby\textquotedblright{}.

\begin{quote}
\fa{مادربزرگ} \fa{هر} \fa{شب} \fa{برای} \fa{نوه} \fa{اش} \fa{لالایی} \fa{میخواند}.
\end{quote}

\begin{quote}
\emph{mâdarbozorg har shab barâ-ye naveash lâlâyi mikhânad.}
\end{quote}

\begin{quote}
Grandmother sings her grandchild a lullaby every night.
\end{quote}

\emph{lâlâyi khândan} — to sing a lullaby. In baby talk \emph{lâlâ kardan} means 'to sleep'.

\begin{cousinweb}
\emph{lâlâyi} is built on \emph{lâlâ}, the soothing sound sung to a baby, with the ending \emph{-i}. It is usually explained as imitative, and EtymDB has no entry for it.
\end{cousinweb}

\begin{grammarlens}[title={What you've built}]
One more word for telling, explaining and describing at B1.
\end{grammarlens}

\subsection*{Your turn}
\begin{itemize}
\raggedright
  \item \emph{Say it:} \emph{lâlâyi}
  \item \emph{Say it:} \emph{lâlâyi}, once more
  \item \emph{Say it:} a sentence of your own with \emph{lâlâyi}
  \item \emph{Recall:} say the Persian for a child, then the Persian for a cradle, then say \emph{lâlâyi} again
\end{itemize}

\begin{tcolorbox}[breakable,colback=teal!4,colframe=teal!35!black,title={Before you move on}]
What does \fa{لالایی} mean? (\textquotedblleft{}A lullaby\textquotedblright{}.) Say it once more.
\end{tcolorbox}
\section[\texorpdfstring{\fa{پوشک} (pushak)}{پوشک (pushak)}]{\fa{پوشک} (pushak) — a nappy, a diaper}

\label{lesson:FA-C502-pushak}



\begin{quote}
Before the new one: say the Persian for a cradle, then the Persian for a lullaby.
\end{quote}

\subsection*{You'll want to know: \fa{پوشک}}
\textbf{\fa{پوشک}} (\emph{pushak}) — \textquotedblleft{}a nappy, a diaper\textquotedblright{}.

\begin{quote}
\fa{پدر} \fa{پوشک} \fa{بچه} \fa{را} \fa{عوض} \fa{کرد}.
\end{quote}

\begin{quote}
\emph{pedar pushak-e bachche râ 'avaz kard.}
\end{quote}

\begin{quote}
The father changed the baby's nappy.
\end{quote}

\emph{pushak 'avaz kardan} — to change a nappy.

\begin{cousinweb}
\emph{pushak} is usually linked with the native verb \emph{pushidan}, 'to cover, to wear', with the ending \emph{-ak}. EtymDB has no entry for it, so the link is hedged.
\end{cousinweb}

\begin{grammarlens}[title={What you've built}]
One more word for telling, explaining and describing at B1.
\end{grammarlens}

\subsection*{Your turn}
\begin{itemize}
\raggedright
  \item \emph{Say it:} \emph{pushak}
  \item \emph{Say it:} \emph{pushak}, once more
  \item \emph{Say it:} a sentence of your own with \emph{pushak}
  \item \emph{Recall:} say the Persian for a cradle, then the Persian for a lullaby, then say \emph{pushak} again
\end{itemize}

\begin{tcolorbox}[breakable,colback=teal!4,colframe=teal!35!black,title={Before you move on}]
What does \fa{پوشک} mean? (\textquotedblleft{}A nappy, a diaper\textquotedblright{}.) Say it once more.
\end{tcolorbox}
\section[\texorpdfstring{\fa{کالسکه} (kâleske)}{کالسکه (kâleske)}]{\fa{کالسکه} (kâleske) — a pram, a pushchair}

\label{lesson:FA-C502-kaleske}



\begin{quote}
Before the new one: say the Persian for a lullaby, then the Persian for a nappy.
\end{quote}

\subsection*{You'll want to know: \fa{کالسکه}}
\textbf{\fa{کالسکه}} (\emph{kâleske}) — \textquotedblleft{}a pram, a pushchair\textquotedblright{}.

\begin{quote}
\fa{او} \fa{بچه} \fa{را} \fa{با} \fa{کالسکه} \fa{به} \fa{پارک} \fa{برد}.
\end{quote}

\begin{quote}
\emph{u bachche râ bâ kâleske be pârk bord.}
\end{quote}

\begin{quote}
She took the baby to the park in the pram.
\end{quote}

Also an old horse-drawn carriage.

\begin{cousinweb}
\emph{kâleske} was borrowed from Russian \emph{koljaska}, 'carriage' (EtymDB).
\end{cousinweb}

\begin{grammarlens}[title={What you've built}]
One more word for telling, explaining and describing at B1.
\end{grammarlens}

\subsection*{Your turn}
\begin{itemize}
\raggedright
  \item \emph{Say it:} \emph{kâleske}
  \item \emph{Say it:} \emph{kâleske}, once more
  \item \emph{Say it:} a sentence of your own with \emph{kâleske}
  \item \emph{Recall:} say the Persian for a lullaby, then the Persian for a nappy, then say \emph{kâleske} again
\end{itemize}

\begin{tcolorbox}[breakable,colback=teal!4,colframe=teal!35!black,title={Before you move on}]
What does \fa{کالسکه} mean? (\textquotedblleft{}A pram, a pushchair\textquotedblright{}.) Say it once more.
\end{tcolorbox}
\section[\texorpdfstring{\fa{شیر} \fa{دادن} (shir dâdan)}{شیر دادن (shir dâdan)}]{\fa{شیر} \fa{دادن} (shir dâdan) — to breastfeed; to feed (milk)}

\label{lesson:FA-C502-shir-dadan}



\begin{quote}
Before the new one: say the Persian for a nappy, then the Persian for a pram.
\end{quote}

\subsection*{You'll want to know: \fa{شیر} \fa{دادن}}
\textbf{\fa{شیر} \fa{دادن}} (\emph{shir dâdan}) — \textquotedblleft{}to breastfeed; to feed (milk)\textquotedblright{}.

\begin{quote}
\fa{مادر} \fa{هر} \fa{سه} \fa{ساعت} \fa{به} \fa{بچه} \fa{شیر} \fa{میدهد}.
\end{quote}

\begin{quote}
\emph{mâdar har se sâ'at be bachche shir midahad.}
\end{quote}

\begin{quote}
The mother feeds the baby every three hours.
\end{quote}

\emph{shir} 'milk' + \emph{dâdan} 'to give'; present stem \emph{deh-}.

\begin{cousinweb}
\emph{shir dâdan} is built on two inherited words: \emph{shir}, 'milk', usually traced to Middle Persian \emph{šīr}, and \emph{dâdan}, 'to give'.
\end{cousinweb}

\begin{grammarlens}[title={What you've built}]
One more word for telling, explaining and describing at B1.
\end{grammarlens}

\subsection*{Your turn}
\begin{itemize}
\raggedright
  \item \emph{Say it:} \emph{shir dâdan}
  \item \emph{Say it:} \emph{shir dâdan}, once more
  \item \emph{Say it:} a sentence of your own with \emph{shir dâdan}
  \item \emph{Recall:} say the Persian for a nappy, then the Persian for a pram, then say \emph{shir dâdan} again
\end{itemize}

\begin{tcolorbox}[breakable,colback=teal!4,colframe=teal!35!black,title={Before you move on}]
What does \fa{شیر} \fa{دادن} mean? (\textquotedblleft{}To breastfeed; to feed (milk)\textquotedblright{}.) Say it once more.
\end{tcolorbox}
